% !TEX program = xelatex
% Spatial Analysis of the Immunosuppressive Microenvironment in Hepatocellular Carcinoma
% Using article document class, compiled with XeLaTeX

\documentclass[12pt, a4paper]{article}

% ============================================================
% Package Loading
% ============================================================
\usepackage{geometry}
\geometry{left=2.5cm, right=2.5cm, top=2.5cm, bottom=2.5cm}

\usepackage{graphicx}
\usepackage{booktabs}
\usepackage{longtable}
\usepackage{amsmath}
\usepackage{amssymb}
\usepackage{hyperref}
\usepackage{xcolor}
\usepackage{listings}
\usepackage{float}
\usepackage{caption}
\usepackage{subcaption}
\usepackage{enumitem}
\usepackage{setspace}
\usepackage{fancyhdr}
\usepackage{natbib}
\usepackage{array}
\usepackage{multirow}
\usepackage{tabularx}
\usepackage{appendix}
\usepackage{microtype}

% ============================================================
% Page and Format Settings
% ============================================================
\onehalfspacing
\pagestyle{fancy}
\fancyhf{}
\fancyhead[L]{\small Spatial Dissection of the Immunosuppressive Microenvironment in HCC via Integrated Single-Cell and Spatial Transcriptomics}
\fancyhead[R]{\small \thepage}
\renewcommand{\headrulewidth}{0.4pt}

\lstset{
    basicstyle=\small\ttfamily,
    keywordstyle=\color{blue},
    commentstyle=\color{gray},
    stringstyle=\color{red},
    breaklines=true,
    frame=single,
    numbers=left,
    numberstyle=\tiny\color{gray},
    language=Python,
    showstringspaces=false,
    xleftmargin=1em,
    framexleftmargin=1em,
    extendedchars=false,
    keepspaces=true
}

\hypersetup{
    colorlinks=true,
    linkcolor=blue,
    citecolor=blue,
    urlcolor=blue
}

% ============================================================
% Title Information
% ============================================================
\title{
    \vspace{-1cm}
    \textbf{Spatial Dissection of the Immunosuppressive Tumour Microenvironment\\
    in Hepatocellular Carcinoma through Integrated Single-Cell\\
    and Spatial Transcriptomics, with Independent Prognostic Validation}\\[0.5em]
    \large Technical Project Report
}
\author{Yifei Tian,\ Nankai University}
\date{}

% ============================================================
% Main Text
% ============================================================
\begin{document}

\maketitle
\thispagestyle{empty}

\begin{abstract}
Hepatocellular carcinoma (HCC) is the third leading cause of cancer-related mortality worldwide, and the objective response rate to immune checkpoint inhibitors remains only approximately 15--20\%. A principal barrier to immunotherapeutic efficacy is the multifaceted immunosuppressive network cooperatively assembled by regulatory T cells (Tregs), tumour-associated macrophages (TAMs) and cancer-associated fibroblasts (CAFs) within the tumour microenvironment (TME). Here we integrate single-cell RNA sequencing (scRNA-seq) with 10x Genomics Visium spatial transcriptomics and employ the Bayesian cell-type deconvolution framework Cell2location to precisely map cell-type annotations onto individual spots of spatial tissue sections. Building on the ``cellular neighbourhood'' framework established by Schürch and colleagues, we apply neighbourhood-composition clustering without reliance on prior-defined thresholds to identify immunosuppressive spatial niches. Quantitative ligand--receptor communication analysis reveals core signalling axes of the Treg--TAM--CAF circuit, including CXCL12--CXCR4, CCL22--CCR4, TGFB1--TGFBR1 and PD-1--PD-L1. Cross-sample reproducibility is validated across two independent Visium sections (HCC4R and CHC20), and the spatial niche signature genes extracted from HCC4R are subsequently projected onto the TCGA-LIHC cohort ($n \approx 370$) via single-sample gene-set enrichment analysis (ssGSEA) for independent prognostic evaluation. Together, this study establishes a complete translational analytical framework spanning spatially resolved discovery, mechanistic communication analysis and clinical-cohort validation.

\noindent\textbf{Keywords:}
hepatocellular carcinoma; spatial transcriptomics; single-cell transcriptomics; immunosuppressive tumour microenvironment; Cell2location;
cellular neighbourhoods; ligand--receptor communication; ssGSEA; survival analysis
\end{abstract}

\newpage
\tableofcontents
\newpage

% ============================================================
%  Part I: Background and Rationale
% ============================================================

\section{Introduction}

\subsection{Disease Background and Clinical Challenges}

Hepatocellular carcinoma (HCC) is the sixth most common malignancy worldwide, with over 900,000 new cases and approximately 830,000 deaths annually; China alone accounts for roughly 45\% of global incident cases~\citep{llovet2021hepatocellular,chen2016liver}. Because the majority of patients present with advanced-stage disease, the five-year overall survival rate remains below 20\%. Sorafenib was formerly the standard first-line treatment for advanced HCC, yet it extended median overall survival by only approximately three months compared with placebo~\citep{llovet2008sorafenib}. The IMbrave150 trial in 2020 demonstrated that atezolizumab combined with bevacizumab raised the objective response rate to 27.3\% and, for the first time, conferred a survival advantage over sorafenib~\citep{finn2020atezolizumab}; nevertheless, approximately 70\% of patients fail to respond to this combination regimen, indicating the existence of robust, multi-layered immunosuppressive mechanisms within the TME.

The immunosuppressive milieu of HCC is orchestrated by multiple cell types acting in concert. Tregs, characterised by the canonical CD$4^+$CD$25^+$FOXP$3^+$ phenotype, directly suppress the anti-tumour activities of effector T cells and NK cells through secretion of IL-10, TGF-$\beta$ and IL-35, as well as expression of checkpoint molecules including CTLA-4, TIGIT and LAG-3~\citep{sakaguchi2020regulatory}. Single-cell transcriptomic studies have demonstrated marked enrichment of Tregs in HCC tissue, and the Treg-to-CD$8^+$ T-cell ratio has been identified as an independent prognostic factor~\citep{zheng2017landscape}. TAMs displaying an M2-polarised phenotype (CD163$^+$MRC1$^+$) actively recruit Tregs into tumour areas via the CCL22--CCR4 and CXCL12--CXCR4 axes, while simultaneously secreting TGF-$\beta$ and upregulating PD-L1 to restrain effector T-cell responses~\citep{noy2014tumor}. The bidirectional positive feedback loop between TAMs and Tregs sustains stable local immunosuppression. CAFs, marked by FAP, ACTA2 and POSTN, recruit immunosuppressive cells via CXCL12 secretion and erect physical barriers against effector T-cell infiltration by remodelling the extracellular matrix (ECM)~\citep{kalluri2016biology,mariathasan2018tgfb}. The cooperative activity of this Treg--TAM--CAF triad represents the core mechanism underlying resistance to immunotherapy in HCC and constitutes the primary biological focus of this study.

\subsection{Spatial Transcriptomics and the Integrative Analysis Paradigm}

Investigating the spatial organisation of immune suppression requires resolving its tissue architecture. Conventional bulk RNA-seq homogenises tissue prior to sequencing, thereby erasing spatial positional information, whereas scRNA-seq, despite its capacity to resolve single-cell transcriptional states, inevitably disrupts spatial relationships during tissue dissociation~\citep{luecken2019current}. The 10x Genomics Visium spatial transcriptomics platform preserves the native tissue architecture and performs whole-transcriptome sequencing at each spatial location (spot, $\sim$55\,$\mu$m diameter), enabling direct observation of cell-type spatial distribution patterns~\citep{williams2022introduction}. Because each spot captures a mixed signal from approximately 1--10 cells, computational deconvolution against an scRNA-seq reference is required to quantitatively estimate cellular composition.

Among current deconvolution tools, Cell2location employs a negative-binomial Bayesian hierarchical model to simultaneously estimate the absolute abundances of multiple cell types across the entire transcriptome, providing posterior uncertainty quantification~\citep{kleshchevnikov2022cell2location}. Multiple independent benchmarks consistently demonstrate its superiority in detecting rare cell types (such as Tregs, typically $<$5\% of the total) over RCTD~\citep{cable2021robust}, CARD~\citep{ma2022integrative} and Tangram~\citep{biancalani2021deep,li2022benchmarking,sang2024spotless,yan2023benchmarking}. The ``cellular neighbourhood'' concept developed by Schürch and colleagues in colorectal cancer employs unsupervised clustering of local cellular composition to reveal how different neighbourhoods contribute differentially to anti-tumour immunity, providing a methodological paradigm for TME subtyping at spatial resolution~\citep{schurch2020coordinated}. Building on this foundation, projecting spatial niche signature genes via ssGSEA~\citep{barbie2009systematic} onto large-scale clinical cohorts and validating their prognostic significance by multivariate Cox regression~\citep{cox1972regression} represents the canonical strategy for translating ``spatial discovery'' into clinical insight~\citep{hanzelmann2013gsva}.

Three critical dimensions remain largely unaddressed in existing HCC studies: elevating the co-localisation of the Treg--TAM--CAF axis from descriptive observation to mechanistic ligand--receptor communication, while preserving spatial information; employing data-driven approaches to identify immunosuppressive spatial niches; and systematically validating spatial transcriptomic findings through independent, large-scale clinical cohort studies. This study is designed to address these unmet needs by establishing a generalisable translational analytical framework.

\subsection{Research Objectives}

This project aims to address four interrelated core scientific questions:
(1) Do immunosuppressive cell signals (Treg, TAM and CAF) exhibit systematic spatial co-localisation within HCC Visium sections, enriching specifically at tumour--stroma interface regions?
(2) Does the Treg--TAM--CAF axis within these regions mediate cooperative immunosuppression through specific ligand--receptor pairs (e.g., CXCL12--CXCR4, CCL22--CCR4)?
(3) Are the immunosuppressive spatial niche features identified in one section reproducible across multiple independent sections?
(4) Can the signature genes derived from spatial data independently predict patient prognosis in the large-scale TCGA-LIHC cohort?

% ============================================================
%  Part II: Methods
% ============================================================

\section{Methods}

\subsection{Overall Analytical Workflow}

We established a complete translational analytical framework encompassing five stages: scRNA-seq reference construction, spatial deconvolution, niche identification, mechanistic dissection and cross-sample/clinical validation (Fig.~\ref{fig:workflow}). The workflow begins with an scRNA-seq reference atlas; Cell2location Bayesian deconvolution maps cell-type abundances onto every spot of the Visium section. Unsupervised clustering of spatial neighbourhood composition then identifies immunosuppressive niches, and a two-layer differential expression strategy extracts niche signature genes. These signature genes subsequently feed three independent but convergent analytical branches: quantitative ligand--receptor communication analysis (mechanistic level), cross-sample reproducibility validation in the independent CHC20 section (spatial level), and ssGSEA projection into the TCGA-LIHC cohort followed by multivariate Cox survival analysis (clinical level). Methodological details for each stage are described in Sections 2.2--2.6 below.

\begin{figure}[h]
    \centering
    \includegraphics[width=1.0\textwidth]{plots/workflow_cn.pdf}
    \caption{Overall analytical workflow.}
    \label{fig:workflow}
\end{figure}

\subsection{Data Sources and Quality Control}

This study integrates single-cell transcriptomic (scRNA-seq), spatial transcriptomic (ST) and large-scale clinical bulk transcriptomic datasets to construct a translational ``peripheral discovery--spatial localisation--clinical validation'' analytical framework.

The single-cell reference dataset was downloaded from the NCBI Gene Expression Omnibus (GEO; accession: GSE149614) and encompasses major cell types including hepatocytes, T/NK cells, B cells, myeloid cells and fibroblasts. Following established best practices for single-cell omics analysis~\citep{luecken2019current}, we rigorously filtered low-quality cells with total expression counts below 200, as well as genes detected in fewer than 10 cells. Because the Cell2location spatial deconvolution model employs a Bayesian negative-binomial regression framework that explicitly models sequencing depth and technical batch effects via probabilistic inference, conventional library-size normalisation or $\log(1+p)$ transformation was not applied at the single-cell pre-processing stage; instead, raw count matrices were retained as direct model inputs.

Spatial transcriptomic data were generated by the 10x Genomics Visium high-throughput platform using a rigorous dual-centre, cross-clinical-context design to establish a closed discovery--validation loop for characterising spatial TME features. The primary discovery section, HCC4R, was derived from a prospective clinical trial cohort (GSE238264) comprising HCC resection specimens obtained following neoadjuvant treatment with cabozantinib combined with nivolumab; this sample is characterised by extensive HCC--CAF spatial interactions. The independent external validation section, CHC20, was sourced from a separate spatial omics cohort focussing on non-alcoholic steatohepatitis-associated hepatocellular carcinoma, with a particular emphasis on myeloid-cell-mediated immunosuppression in tumour fibrotic foci.

Following established best practices and quality-control standards for spatial transcriptomics~\citep{williams2022introduction,palla2022squidpy}, both ST sections were uniformly filtered: only spots with total UMI counts between 5,000 and 35,000 and mitochondrial gene fraction below 20\% were retained, and genes detected in fewer than 10 spots were excluded.

For independent large-scale clinical validation of the spatial omics findings, we incorporated the TCGA-LIHC cohort (approximately 370 HCC patients). Standardised HTSeq-FPKM expression matrices and complete clinical follow-up and survival data were downloaded from the Genomic Data Commons (GDC) via the TCGAbiolinks R package~\citep{colaprico2016tcgabiolinks}.

\subsection{Treg Subtype Identification and Cell2location Deconvolution}

T/NK cell subsets were extracted from scRNA-seq data and subjected to high-resolution subclustering using the Leiden algorithm (resolution = 3.0)~\citep{traag2019leiden}. The expression distributions of \textit{CD3D}, \textit{CD4}, \textit{FOXP3} and \textit{IL2RA} across individual subclusters were visualised using bubble plots. Upon confirmation that a specific cluster exhibited the canonical Treg phenotype CD4$^+$CD25$^+$FOXP3$^+$~\citep{sakaguchi2020regulatory}, the cluster was re-annotated as Treg and incorporated into the final cell-type labels.

Cell2location was executed in two stages~\citep{kleshchevnikov2022cell2location}. In the first stage, a RegressionModel was trained on raw scRNA-seq count data to derive a cell-type reference signature matrix, using a maximum of 250 training epochs, a batch size of 1,024 and early stopping with patience = 30. In the second stage, the reference signature matrix was used for independent spatial modelling of each Visium section, with the key prior parameter N\_cells\_per\_location set to 30 (informed by the Visium spot diameter of $\sim$55\,$\mu$m and typical hepatic tissue cell density~\citep{kleshchevnikov2022cell2location}), detection\_alpha set to 20 (the recommended default), and a maximum of 1,000 training epochs. The posterior mean abundance matrix (\texttt{means\_cell\_abundance\_w\_sf}) was used as the estimate of absolute cell-type abundance per spot.

To facilitate biological interpretation and cross-method cross-validation, we additionally assigned each spot a supplementary rule-based spatial domain label (tumor\_core / tumor\_edge / stroma\_immune / other), beyond the primary Cellular Neighbourhoods clustering. This annotation was used solely for visualisation and cross-validation purposes and was not included in the main niche-discovery pipeline. Adapting the local cellular composition threshold approach used by Schürch et al.~\citep{schurch2020coordinated} and Bhate et al.~\citep{bhate2022tissue} to demarcate tumour parenchyma, invasive front and stroma/immune regions in multiplexed spatial omics studies: spots with hepatocyte proportions $\geq$ the 75th percentile were designated \textit{hep\_high} and labelled \textit{tumor\_core}; among the remaining spots, those with an immune\_stroma\_score (see Appendix~\ref{app:scores}) $\geq$ the 50th percentile were labelled \textit{stroma\_immune}; spots satisfying this criterion and having at least one \textit{hep\_high} neighbour within the physical neighbourhood radius were re-classified as \textit{tumor\_edge}, to delineate the transitional interface between tumour parenchyma and immune stroma; remaining spots were labelled \textit{other}.

\subsection{Identification of Immunosuppressive Spatial Niches}

Cell2location-derived absolute abundances were normalised to proportional estimates, after which a $k$-nearest-neighbour ($k$-NN) graph was constructed using tissue spatial coordinates. For each spot $i$, the mean cell-type proportion across its $k$ physically nearest neighbours was computed to yield a neighbourhood composition vector $\mathbf{n}_i$:
\begin{equation}
\mathbf{n}_i = \frac{1}{|\mathcal{N}_i|} \sum_{j \in \mathcal{N}_i} \mathbf{p}_j
\label{eq:neighborhood_vector}
\end{equation}
Following the classic parameter choice of Schürch et al.~\citep{schurch2020coordinated} in colorectal cancer, we set $k = 15$. A new $k$-NN graph was then constructed in the high-dimensional feature space defined by $\mathbf{n}_i$, and Leiden community detection (resolution $= 0.5$) was applied to this graph to adaptively identify spatial niche types. Compared with $k$-means, Leiden does not require the number of clusters to be specified a priori and guarantees cluster connectivity by construction~\citep{traag2019leiden}.

For each Leiden cluster, the mean values of four functional dimensions---Treg proportion, myeloid proportion, fibroblast proportion and immunosuppressive gene module score---were computed. The cluster with the minimum sum of descending ranks across all four dimensions was designated the immunosuppressive niche:
\begin{equation}
R_c = \text{rank}(\bar{p}_{\text{Treg},c}) + \text{rank}(\bar{p}_{\text{Myeloid},c})
+ \text{rank}(\bar{p}_{\text{Fibroblast},c}) + \text{rank}(\bar{s}_{\text{immune},c})
\label{eq:combined_rank}
\end{equation}
This multi-dimensional composite score ensures that the identified niche simultaneously satisfies the criteria of co-enrichment across the Treg--TAM--CAF triad, rather than reflecting isolated overexpression of a single cell type.

The immunosuppressive gene module score was defined as the CP10K+log1p mean expression of 14 literature-supported marker genes (\textit{FOXP3}, \textit{IL2RA}, \textit{CTLA4}, \textit{TIGIT}, \textit{LAG3}, \textit{PDCD1}, \textit{HAVCR2}, \textit{TGFB1}, \textit{IL10}, \textit{CXCL12}, \textit{CCL22}, \textit{TNFRSF18}, \textit{TNFRSF4}, \textit{IKZF2})~\citep{sakaguchi2020regulatory,sharma2023immunosuppressive}. Auxiliary scores (Treg-like score, immune-stroma score and niche score) were computed by summing Z-score-normalised values of the corresponding compositional dimensions (see Appendix).

To verify that the neighbourhood parameter choices did not dominate the conclusions, we conducted a continuous-valued sensitivity analysis free of hard-threshold cutoffs. The composite niche score vector under the primary analysis parameters ($k = 15$, using the four-dimensional Z-score sum underlying Equation~\ref{eq:combined_rank}) was used as the reference baseline. Perturbations were applied across $k \in \{10, 15, 20\}$ and radius multipliers $\in \{1.0, 1.25, 1.5\}$, yielding six parameter combinations. Consistency between each perturbed score distribution and the reference was jointly quantified using Spearman's rank correlation coefficient $\rho$ and a continuous weighted Jaccard similarity WJ (formulae defined in Appendix~\ref{app:sensitivity}). Robustness criteria were as follows: $\rho > 0.95$ and $\text{WJ} > 0.90$ indicates highly robust results; $\rho > 0.90$ and $\text{WJ} > 0.80$ indicates robust results insensitive to parameter choice; $\rho < 0.85$ or $\text{WJ} < 0.70$ indicates a material impact of the parameter change warranting re-examination. The threshold $\rho > 0.90$ and $\text{WJ} > 0.80$ was adopted as the passing criterion for parameter robustness.

Following neighbourhood parameter robustness validation, a binary high/low niche label was assigned to each spot based on the continuous composite immunosuppressive niche score (\textit{immunosuppressive\_niche\_score}), serving as the grouping variable for subsequent two-layer signature gene selection (Section 2.6). Specifically, spots with scores $\geq$ the 90th percentile across all spots were labelled \textit{niche\_high}; remaining spots were labelled \textit{niche\_low}.

\subsection{Two-Layer Signature Gene Selection}

Each Visium spot captures the mixed signal of approximately 5--50 cells (cell dilution effect), resulting in near-zero raw expression of Treg markers such as \textit{FOXP3} and \textit{CTLA4} in the vast majority of spots. This severely attenuates the mean difference ($\log_2$FC) between groups and prevents traditional strategies based solely on Wilcoxon testing and $\log_2$FC thresholds from capturing these low-abundance yet biologically critical genes within the top-$N$ candidate list. Accordingly, beyond $\log_2$FC, we incorporated the detection rate (fraction of spots with expression $> 0$) as a complementary dimension, and adopted a two-layer filtering strategy to extract immunosuppressive niche signature genes.

\textbf{Layer 1 (Integrated Wilcoxon + BH-FDR + AUC + detection-rate ranking):}
Geographically adjacent spots were first aggregated into block-level means via spatial grid aggregation (block\_mwu mode) to reduce the violation of hypothesis-testing independence caused by spatial autocorrelation~\citep{palla2022squidpy}. Blocks were classified as \textit{niche\_high} (score $\geq P_{90}$) or \textit{niche\_low} (score $< P_{90}$) based on the 90th percentile of the composite immunosuppressive niche score. A one-sided Mann--Whitney U test (i.e., Wilcoxon rank-sum test) was applied to each gene, and $p$-values were corrected for multiple testing using the Benjamini--Hochberg (BH) procedure~\citep{mann1947test}. Detection rates $\text{frac}_{\text{high}}$ and $\text{frac}_{\text{low}}$ and their difference $\Delta\text{frac} = \text{frac}_{\text{high}} - \text{frac}_{\text{low}}$ were computed, together with the AUC derived by standardising the Mann--Whitney U statistic~\citep{hanley1982meaning}. Primary inclusion criteria were FDR $< 0.1$ and AUC $> 0.60$ and ($\log_2$FC $> 0.3$ OR $\Delta$frac $> 0.10$); if fewer than 80 genes met these criteria, they were relaxed to FDR $< 0.2$ and AUC $> 0.55$ and ($\log_2$FC $> 0.2$ OR $\Delta$frac $> 0.05$). Candidate genes were further filtered by removing a blacklist of over 350 non-specific genes encompassing housekeeping genes, ribosomal proteins, mitochondrial genes, haemoglobin genes, cell-cycle genes, broadly expressed inflammatory genes and tissue-damage/stress-response genes, to ensure niche specificity rather than ubiquitous tissue activity. Candidates were ranked by a weighted composite score:
\begin{equation*}
\text{signature\_rank\_score} = 0.35 \cdot \text{minmax}(\log_2\text{FC}) + 0.35 \cdot \text{minmax}(\Delta\text{frac}) + 0.30 \cdot \text{minmax}(\text{AUC})
\label{eq:signature_rank_score}
\end{equation*}
The top 80 genes (without truncation if forced inclusion of prior genes pushes the total above 80) constituted the preliminary signature gene set.

\textbf{Layer 2 (Targeted AUC testing and mandatory inclusion of prior functional genes):}
For three groups of functionally annotated prior gene sets---Treg markers (\textit{FOXP3}, \textit{IL2RA}, \textit{CTLA4}, etc.), TAM signature genes (\textit{CD163}, \textit{MRC1}, \textit{TGFB1}, etc.) and CAF activation genes (\textit{FAP}, \textit{ACTA2}, \textit{COL1A1}, etc.)---AUC, $\log_2$FC, $\Delta$frac and BH-corrected $p$-values were independently computed between \textit{niche\_high} and \textit{niche\_low} groups to assess niche-specific enrichment. Prior genes satisfying AUC $> 0.60$ and $\Delta$frac $> 0.05$ (and not on the non-specific gene blacklist) were mandatorily incorporated into the final signature gene list regardless of their Layer 1 genome-wide ranking, preventing the omission of biologically significant key genes due to Visium's sparse expression profile. Complete implementation of the two-layer filtering pipeline is provided in Appendix~\ref{app:code}.

\subsection{Ligand--Receptor Communication Analysis}

To elevate the descriptive co-localisation findings to the level of mechanistic intercellular communication, we quantitatively compared the communication intensities of 12 literature-supported ligand--receptor (L-R) pairs---including \textit{CXCL12}--\textit{CXCR4}, \textit{CCL22}--\textit{CCR4}, \textit{TGFB1}--\textit{TGFBR1} and PD-1--PD-L1 (complete list in Appendix~\ref{app:lr_table})---between \textit{niche\_high} and \textit{niche\_low} groups. The conventional ligand--receptor product method~\citep{efremova2020cellphonedb} is prone to producing products of zero when either gene is near-zero at the Visium spot level (due to cell dilution effects), effectively annihilating signals from Treg-specific pathways such as \textit{CCL22}--\textit{CCR4}. We therefore rank-normalised each gene's expression vector across all spots (mapping values to $[0,1]$) and used the product of ligand and receptor ranks as a proxy communication score:
\begin{equation*}
\text{score}_{LR,i} = \text{rank}(E_{L,i}) \times \text{rank}(E_{R,i}), \quad \text{rank}(\cdot) \in [0,1]
\end{equation*}
This rank-normalised product preserves the relative ordering of ligand and receptor expression, effectively circumvents the dominance of sparse zeros on multiplicative scores, and naturally enforces the logical AND constraint that both ligand and receptor must be simultaneously highly expressed. For each L-R pair, the mean rank-normalised product was computed in \textit{niche\_high} and \textit{niche\_low} groups, and the enrichment or depletion direction in the immunosuppressive niche was quantified as $\log_2(\text{mean}_{\text{high}}/\text{mean}_{\text{low}})$ (with a pseudocount $\varepsilon = 10^{-6}$ added to the denominator to avoid division by zero).

Beyond intercellular communication, we systematically characterised pathway-level biological process activities---including immunosuppression, stromal remodelling, inflammation and metabolic stress---across spatial domains using pre-defined functional gene-set scores (Fig.~\ref{fig:domain_func}). Guided by the construction principles of Hallmark gene sets~\citep{liberzon2015molecular} and the general paradigm of single-sample gene-set scoring~\citep{barbie2009systematic}, we curated 10 functional pathway gene sets closely linked to the Treg--TAM--CAF immunosuppressive axis and its associated microenvironmental alterations: TGF-$\beta$ Signalling, EMT, Angiogenesis, IL6-JAK-STAT3, Treg Immune Suppression, TAM Polarisation (M2), CAF Activation, Immune Checkpoint, Hypoxia and TNF-$\alpha$ Signalling via NF-$\kappa$B, each comprising six representative genes. For each spot, the CP10K+log1p expression mean of all available pathway genes was computed as the single-sample score for that pathway, then averaged across spatial domain groups (tumor\_core / tumor\_edge / stroma\_immune / other) as defined in Section 2.3. In bubble plots, bubble size represents the relative magnitude of the score, and colour indicates whether the score exceeds the global mean across all four domains for that pathway (red: above average; blue: below average).

\subsection{Cross-Section Consistency Validation and TCGA Survival Analysis}

Using the signature gene set extracted from HCC4R as a reference, immune module genes (Treg/TAM-related, $n = 21$) and stromal/ECM module genes (CAF/ECM-related, $n = 21$) were scored separately on the CHC20 section. Z-score-standardised values of the two module scores were summed to yield a combined niche score (immune\_stromal\_niche\_score). Cross-sample consistency of score distributions was quantified via spatial projection maps and Mann--Whitney U tests~\citep{mann1947test}.

Survival analysis was conducted using the ssGSEA method~\citep{barbie2009systematic} implemented in the GSVA R package~\citep{hanzelmann2013gsva}, with $\log_2(\text{FPKM}+1)$ expression matrices as input, a Gaussian kernel (kcdf = Gaussian) and absolute ranking (abs.ranking = TRUE), to compute a scalar enrichment score for each TCGA-LIHC patient. ssGSEA generates an independent continuous score reflecting niche activity per patient, making it more appropriate for incorporation into prognostic regression models than the cross-sample relative activity estimates produced by GSVA~\citep{barbie2009systematic}. A multivariate Cox proportional hazards model~\citep{cox1972regression} was constructed with ssGSEA score, age and tumour stage as covariates:
\begin{equation*}
h(t \mid \mathbf{x}) = h_0(t) \cdot \exp(\beta_1 s_{\text{ssGSEA}} + \beta_2 \text{age} + \beta_3 \text{stage})
\label{eq:cox}
\end{equation*}
After controlling for age and tumour stage (established strong independent prognostic factors in HCC~\citep{llovet2021hepatocellular}), a statistically significant $\beta_1$ for the ssGSEA score would indicate that the immunosuppressive niche signature captures prognostic information independent of clinical staging. Kaplan--Meier survival curves~\citep{kaplan1958nonparametric} were plotted by dichotomising patients at the ssGSEA score median, and the log-rank test~\citep{mantel1966evaluation} was used to assess survival differences between the two groups.

% ============================================================
%  Part III: Results
% ============================================================

\section{Results}

\subsection{Construction of the scRNA-seq Reference Atlas and Identification of Treg Subsets}

\begin{figure}[h]
    \centering
    \begin{subfigure}[b]{0.48\textwidth}
        \centering
        \includegraphics[width=\textwidth]{plots/cell_type_marker_gene_expression.png}
        \caption{Heatmap of canonical marker gene expression across major cell types.}
        \label{fig:cell_type}
    \end{subfigure}
    \hfill
    \begin{subfigure}[b]{0.48\textwidth}
        \centering
        \includegraphics[width=\textwidth]{plots/treg_marker_gene_expression_t_nk_subclusters.png}
        \caption{Bubble plot of Treg marker gene expression in T/NK cell subclusters.}
        \label{fig:treg}
    \end{subfigure}
    \caption{Cell-type annotation of the single-cell reference atlas and high-resolution identification of the regulatory T-cell (Treg) subpopulation.}
    \label{fig:scRNA-seq_Treg}
\end{figure}

A cell-type reference atlas was constructed from the GSE149614 single-cell transcriptomic dataset. The cell-type marker gene heatmap (Fig.~\ref{fig:cell_type}) displays Z-score-normalised mean expression profiles for the seven major cell types resolved (including T/NK, Treg and myeloid cells), with canonical marker genes on the vertical axis. Hepatocytes/malignant cells (\textit{ALB}, \textit{APOA2}, \textit{APOA1}), T/NK cells (\textit{CD3D}, \textit{GNLY}, \textit{GZMB}), myeloid cells (\textit{LYZ}, \textit{AIF1}, \textit{HLA-DRA}), HSC/fibroblasts (\textit{COL1A1}, \textit{ACTA2}, \textit{RGS5}), B cells (\textit{MS4A1}, \textit{CD79A}) and endothelial cells (\textit{PECAM1}, \textit{CDH5}) all display clear diagonal high-expression patterns, confirming the reliability of cell-type annotations~\citep{ma2019comprehensive}.

For the T/NK subset, high-resolution Leiden clustering at resolution = 3.0 was performed and the results visualised as a horizontal bubble plot (Fig.~\ref{fig:treg}), showing the expression of \textit{CD3D}, \textit{CD4}, \textit{FOXP3} and \textit{IL2RA} across T/NK subclusters. Bubble size represents the fraction of expressing cells and colour intensity represents mean expression level. A specific cluster (cluster 9) exhibited the largest and darkest \textit{FOXP3} bubble alongside co-expression of \textit{CD4} and markedly elevated \textit{IL2RA} relative to all other clusters, consistent with the canonical CD4$^+$CD25$^+$FOXP3$^+$ Treg phenotype~\citep{sakaguchi2020regulatory}. \textit{FOXP3} expression was near-absent in all other T/NK subclusters, demonstrating the subpopulation specificity of \textit{FOXP3} expression. Cluster 9 was re-annotated as Treg and the final cell-type labels were incorporated into the Cell2location deconvolution reference.

\subsection{Cell2location Deconvolution Reveals the Spatial Distribution of Immune Cells in HCC4R}
\subsubsection{Spatial Distribution of Immune Cell Types in HCC4R}

\begin{figure}[h]
    \centering
    \includegraphics[width=0.6\textwidth]{plots/hcc4r_he_image.png}
    \caption{Anatomical demarcation of spatial niches and tumour microenvironment domains in HCC4R.}
    \label{fig:he}
\end{figure}

Figure~\ref{fig:he} precisely maps the spatial niche landscape of the HCC4R section onto individual spots. The H\&E-stained image clearly delineates the morphological boundary between the tumour parenchyma (large, darkly stained nuclei; dense hepatocyte packing) and the stroma/interstitium (fibrosis; inflammatory infiltration), providing a histological reference for the spatial transcriptomic results.

The spatial domain map reveals that the tumour core (red) occupies a large, contiguous area forming the principal scaffold of the tissue section, spatially coinciding with the nucleus-dense region evident in H\&E staining and hepatocyte deconvolution abundance maps. The stroma/immune domain (blue) is embedded in bands surrounding the tumour core and at tissue edges, forming a continuous zone that encircles and interpenetrates the tumour parenchyma rather than appearing as isolated scattered spots; this pattern is consistent with the expectation that unsupervised clustering effectively captures continuous spatial microstructure rather than random noise~\citep{schurch2020coordinated}. The tumour edge domain (orange) is scattered at the transitional zone between the tumour core and stroma/immune domain, forming a narrow interface band that represents the smallest proportion in the section. This spatial pattern suggests that the tissue architecture of HCC4R does not exhibit a binary tumour--stroma partition, but instead a continuous ``core--edge--stroma-immune'' gradient, consistent with the classical ``immune-excluded'' HCC spatial organisation in which tumour cells occupy the core while immune-stromal cells are confined to the periphery~\citep{llovet2021hepatocellular,binnewies2018understanding}.

\begin{figure}[h]
    \centering
    \begin{subfigure}[b]{0.48\textwidth}
        \centering
        \includegraphics[width=\textwidth]{plots/spatial_hepatocyte.png}
        \caption{Spatial proportion distribution of hepatocytes.}
        \label{fig:spatial_hepatocyte}
    \end{subfigure}
    \hfill
    \begin{subfigure}[b]{0.48\textwidth}
        \centering
        \includegraphics[width=\textwidth]{plots/spatial_treg.png}
        \caption{Spatial abundance distribution of Treg-like cells.}
        \label{fig:spatial_treg}
    \end{subfigure}

    \vspace{0.3cm}

    \begin{subfigure}[b]{0.48\textwidth}
        \centering
        \includegraphics[width=\textwidth]{plots/spatial_myeloid.png}
        \caption{Spatial proportion distribution of myeloid cells.}
        \label{fig:spatial_myeloid}
    \end{subfigure}
    \hfill
    \begin{subfigure}[b]{0.48\textwidth}
        \centering
        \includegraphics[width=\textwidth]{plots/spatial_fibroblast.png}
        \caption{Spatial proportion distribution of fibroblasts.}
        \label{fig:spatial_fibroblast}
    \end{subfigure}
    \caption{Spatial abundance and proportion maps of major cell types resolved by Bayesian deconvolution.}
    \label{fig:spatial_cell_abundance}
\end{figure}

The cell-type spatial distribution maps (Fig.~\ref{fig:spatial_cell_abundance}) reveal spatial patterns in close agreement with biological expectations. Hepatocyte high-abundance regions are concentrated in the tumour parenchyma, spatially coinciding with the nucleus-dense area in H\&E staining (Fig.~\ref{fig:spatial_hepatocyte}). Tregs display focal enrichment (right-skewed distribution), with high-abundance hotspots primarily at the tumour-parenchyma/stroma interface (Fig.~\ref{fig:spatial_treg}), consistent with the biological propensity of Tregs to accumulate at the ``tumour--immune interface''~\citep{zheng2017landscape}. Myeloid cells and Tregs exhibit marked spatial co-localisation (Fig.~\ref{fig:spatial_myeloid}), echoing the mechanistic role of TAMs in actively recruiting Tregs via the CCL22--CCR4 axis~\citep{noy2014tumor}. Fibroblasts are predominantly distributed within the stromal layer surrounding the tumour parenchyma (Fig.~\ref{fig:spatial_fibroblast}), reflecting the functional positioning of CAFs in constructing physical barriers~\citep{kalluri2016biology}.

\begin{figure}[h]
    \centering
    \includegraphics[width=0.8\textwidth]{plots/cell_type_spatial_colocalization.png}
    \caption{Spatial co-localisation correlation heatmap of major cell types in the tissue section.}
    \label{fig:hot_plot}
\end{figure}

Pearson correlation analysis of cell-type pairs (Fig.~\ref{fig:hot_plot}) further quantifies the spatial co-localisation relationships described above. The correlation heatmap reveals that myeloid cells exhibit moderate positive correlations with both Tregs ($r = 0.35$) and fibroblasts ($r = 0.35$), while fibroblasts and Tregs display a weak negative correlation ($r = -0.15$). This ``asymmetric triangle'' correlation pattern suggests fine-grained functional stratification of immune-stromal components in space: myeloid cells (TAMs) function as a central spatial hub simultaneously associated with immunosuppressive and stromal cells---on the one hand recruiting and anchoring Tregs via secretion of CCL22, CXCL12 and related chemokines~\citep{noy2014tumor,curiel2004specific}, and on the other cooperating with fibroblasts (CAFs) in ECM remodelling and physical barrier construction~\citep{kalluri2016biology}---while fibroblasts and Tregs have a weaker direct spatial association, suggesting that they are more likely indirectly coupled through TAMs as a shared intermediary node rather than directly recruiting one another. This observation is consistent with reports of TAMs serving as a ``signalling hub'' in the tumour immunosuppressive network.

Furthermore, all three immune-stromal components (myeloid, Treg and fibroblast cells) display strong negative spatial correlations with hepatocytes ($r = -0.73$, $-0.62$ and $-0.55$, respectively), indicating that these three populations are systematically concentrated in non-parenchymal regions and excluded from the tumour core. This quantitative evidence strongly supports the conclusion that the immunosuppressive niche predominantly occupies the tumour--stroma transitional zone rather than being diffusely distributed throughout the section---a topological feature that impedes the effective infiltration of effector immune cells into the tumour core, consistent with the classical description of ``immune-excluded'' tumours in which immune cells are physically restricted to the tumour boundary and stroma, unable to enter the parenchyma to exert cytotoxic functions~\citep{llovet2021hepatocellular}.

\subsubsection{Differential Characteristics Across Spatial Domains in HCC4R}

\begin{figure}[h]
    \centering
    \begin{subfigure}[b]{0.8\textwidth}
        \centering
        \includegraphics[width=\textwidth]{plots/cell_type_composition_per_spatial_domain.png}
        \caption{Cell-type composition per spatial domain.}
        \label{fig:cell_type_composition}
    \end{subfigure}
    \vspace{0.5cm}

    \begin{subfigure}[b]{0.32\textwidth}
        \centering
        \includegraphics[width=\textwidth]{plots/treg_like_score_by_spatial_domain.png}
        \caption{Treg-like score.}
        \label{fig:treg_score}
    \end{subfigure}
    \hfill
    \begin{subfigure}[b]{0.32\textwidth}
        \centering
        \includegraphics[width=\textwidth]{plots/immune_stroma_score_by_spatial_domain.png}
        \caption{Immune-stroma score.}
        \label{fig:immune_stroma}
    \end{subfigure}
    \hfill
    \begin{subfigure}[b]{0.32\textwidth}
        \centering
        \includegraphics[width=\textwidth]{plots/immunosuppressive_niche_score_by_spatial_domain.png}
        \caption{Immunosuppressive niche score.}
        \label{fig:immunosuppressive_niche}
    \end{subfigure}
    \caption{Cell-type composition and module scores across spatial domains.}
    \label{fig:spatial_domain_scores}
\end{figure}

The stacked bar chart of cell-type composition (Fig.~\ref{fig:cell_type_composition}) provides quantitative support for the spatial patterns described above. The hepatocyte proportion exhibits a clear monotonic decrease across the four domains: tumour core ($\sim$35\%) > other ($\sim$24\%) > tumour edge ($\sim$25\%) > stroma/immune (lowest), consistent with the definition of the tumour core as being dominated by hepatocytes/malignant cells. Conversely, the proportions of all three immunosuppressive axis cell types---Tregs, myeloid cells and fibroblasts---are higher in the stroma/immune and tumour edge domains than in the tumour core. Treg proportions reach their highest level in the stroma/immune domain, with myeloid cells and fibroblasts showing similar enrichment trends; the cooperative enrichment of these three populations within the same spatial domain directly corroborates the spatial synergy of the Treg--TAM--CAF immunosuppressive triad described in the Introduction~\citep{sakaguchi2020regulatory,noy2014tumor,kalluri2016biology}. T/NK cell proportions are relatively elevated in the tumour edge, consistent with the view that the invasive front is the principal battleground between effector and regulatory immune cells~\citep{schurch2020coordinated}. Endothelial cell proportions are comparatively stable across the four domains, suggesting that vascular distribution per se is not the dominant factor driving the clustering partition.

Violin plots of module scores across spatial domains (Figs.~\ref{fig:treg_score}, \ref{fig:immune_stroma}, \ref{fig:immunosuppressive_niche}) yield analogous conclusions. Kruskal--Wallis testing reveals highly significant overall distributional differences among domains for all three scores ($p < 0.001$ in each case). The stroma/immune domain exhibits significantly higher immune-stroma scores ($\eta^2 = 0.83$), immunosuppressive niche scores ($\eta^2 = 0.41$) and Treg-like scores ($\eta^2 = 0.33$) than all other domains (pairwise $p < 0.001$). This disparity is most pronounced for the immune-stroma score: non-stroma/immune spots are highly concentrated in the low-score range, while stroma/immune spots show a pronounced rightward shift, spanning the high-score range~\citep{mann1947test}.

Integrating the quantitative evidence from both cell-type composition and functional scores, HCC4R displays a monotonic immunosuppressive gradient of ``low tumour core, intermediate tumour edge, high stroma/immune'', supported by statistically significant effect sizes. This provides prior expectation for the data-driven Cellular Neighbourhoods clustering presented in the next section.

\subsection{Cellular Neighbourhoods Clustering Identifies the Immunosuppressive Spatial Niche}

\begin{figure}[h]
    \centering
    \includegraphics[width=0.8\textwidth]{plots/neighborhood_composition_clusters.png}
    \caption{Leiden clustering of immunosuppressive spatial niches based on neighbourhood composition.}
    \label{fig:niche_cluster}
\end{figure}

Figure~\ref{fig:niche_cluster} displays the clustering results obtained by running Leiden community detection (resolution $= 0.5$) on the $k$-NN graph reconstructed in the feature space of neighbourhood composition vectors $\mathbf{n}_i$ ($k = 15$) described in Section 2.4. Spots are coloured by cluster index (NC-0 to NC-11; 12 neighbourhood clusters in total) in spatial coordinates, fully preserving the heterogeneous tissue structure discovered natively by the Leiden algorithm in the neighbourhood composition feature space.

The most prominent feature of the figure is the immunosuppressive niche, which is the single spatially largest and most highly connected cluster, occupying the central area of the section and extending in multiple directions. It surpasses all other 11 clusters in spatial continuity and coverage. This cluster was determined by the composite ranking formula (Equation~\ref{eq:combined_rank}) to be the \textit{immunosuppressive\_niche}---the cluster with the minimum sum of descending ranks across four dimensions (Treg proportion, myeloid proportion, fibroblast proportion and immunosuppressive gene score)---and thus the neighbourhood type with the strongest overall immunosuppressive characteristics under joint four-dimensional evidence. Spatially, the immunosuppressive niche is contained within the stroma/immune and tumour edge domains in Fig.~\ref{fig:he}. The spatial concordance between these two independent annotation approaches (raw Leiden cluster indices vs. rule-based semantic labels) methodologically reinforces the robustness and interpretability of the immunosuppressive niche determination, and aligns with the expectation that clustering and semantic annotation maps should exhibit high spatial agreement~\citep{schurch2020coordinated,goltsev2018deep}.

A further notable feature of the figure is that, aside from the red-marked immunosuppressive niche, the remaining clusters display clear histological zoning: NC-1 and NC-3, for example, form two large, contiguous independent areas in the lower-left and upper-right corners of the section, likely corresponding to distinct tumour parenchyma subtypes from different regions or batches; NC-5 and NC-9 are embedded as narrow bands between larger clusters such as NC-1 and the immunosuppressive niche, suggesting these neighbourhoods represent transitional sub-structures between tumour and stroma. Overall, the Leiden algorithm autonomously identified 12 internally coherent and spatially contiguous neighbourhood types without any prior specification of cluster numbers, validating that the unsupervised clustering approach genuinely captures real tissue spatial microstructures rather than artificially imposed artefacts~\citep{traag2019leiden,schurch2020coordinated}. The co-existence of ``one prominent immunosuppressive neighbourhood and multiple tumour/stromal sub-structure neighbourhoods'' is consistent with recent spatial proteomics observations that tumour tissues universally harbour multiple co-existing cellular neighbourhoods, and that immunosuppressive neighbourhoods typically form continuous structures spanning the tumour--stroma interface rather than isolated hotspots~\citep{jackson2020single,bhate2022tissue}, further supporting the conclusion that the immunosuppressive niche identified here is a genuine, histologically grounded spatial entity in HCC4R and not an artefact of arbitrary clustering parameter choices.

\subsection{Extraction and Functional Annotation of Niche Signature Genes}

\begin{figure}[h]
    \centering
    \begin{subfigure}[b]{0.71\textwidth}
        \centering
        \includegraphics[width=\textwidth]{plots/fraction_of_spots_enrichment.png}
        \caption{Scatter plot of signature gene selection between high and low niche groups.}
        \label{fig:high_low_niche_gene}
    \end{subfigure}
    \hfill
    \begin{subfigure}[b]{0.27\textwidth}
        \centering
        \includegraphics[width=\textwidth]{plots/niche_signature_genes_by_spatial_domain.png}
        \caption{Heatmap of top differentially expressed genes across spatial domains.}
        \label{fig:domain_hotplot}
    \end{subfigure}
    \caption{Identification and landscape of differentially expressed genes across tumour microenvironment sub-regions.}
    \label{fig:niche_core_gene}
\end{figure}

Figure~\ref{fig:high_low_niche_gene} plots $\log_2$FC on the $x$-axis against $\Delta$fraction (detection rate difference) on the $y$-axis, depicting the Layer 1 detection-rate-compensated filtering results. Red dots (both-significant, $n = 33$) satisfy both significance criteria simultaneously, including high-confidence immunosuppressive genes such as \textit{CXCL12} and \textit{TGFB1}, consistent with known chemokine axes and CAF/Treg secretory mechanisms~\citep{domanska2013targeting,mariathasan2018tgfb}. Notably, the $\log_2$FC values of most immunosuppressive candidate genes are concentrated in the narrow range of 0.1--0.3 despite relatively high $\Delta$fraction values; this counterintuitive pattern is explained by the Visium cell dilution effect~\citep{luecken2019current}, underscoring the necessity of incorporating the $\Delta$frac dimension as a complementary correction to conventional $\log_2$FC-based filtering.

Figure~\ref{fig:domain_hotplot} presents a heatmap of Z-score expression for the top signature genes (ranked in descending order of composite score) across the four spatial domains. Nearly all signature genes consistently exhibit the highest expression in the stroma/immune domain and the lowest in the tumour core, with intermediate values in the tumour edge, in close agreement with the spatial landscape in Fig.~\ref{fig:he}~\citep{zheng2017landscape}. The heatmap also reveals co-enrichment of MHC class II molecules, lymphocyte-associated genes (\textit{CD52}, \textit{TRAC} and others) and stromal/fibroblast genes (\textit{DCN}, \textit{LUM} and others) in the stroma/immune domain, suggesting that this region combines active antigen presentation with stromal remodelling~\citep{neefjes2011towards,iozzo2015proteoglycan}.

\begin{figure}[h]
    \centering
    \includegraphics[width=0.9\textwidth]{plots/pathway_enrichment_by_spatial_domain.png}
    \caption{Enrichment landscape of pre-defined functional pathway gene-set scores across spatial domains.}
    \label{fig:domain_func}
\end{figure}

Figure~\ref{fig:domain_func} presents a bubble plot of enrichment scores for 10 pre-defined functional pathway gene sets across the four spatial domains. Globally, the tumour core exhibits the lowest activity across all 10 pathways, while 9 of 10 pathways (all except Hypoxia) peak in the stroma/immune domain, confirming this region as the most active site for pro-tumorigenic, stromal-remodelling and immunosuppressive signalling. In terms of differential patterns, CAF Activation is exclusively enriched in the stroma/immune domain, consistent with CAF secretion of CXCL12 and ECM remodelling~\citep{kalluri2016biology}; EMT, TGF-$\beta$ and Treg Immune Suppression pathways are elevated in both the tumour edge and stroma/immune domains, consistent with the expected cooperative activity of the invasive front and immunosuppressive core~\citep{nieto2016emt,sakaguchi2020regulatory}. Notably, the Hypoxia pathway is enriched in the tumour edge rather than the stroma/immune domain, possibly reflecting more acute HIF pathway activation at sites of active angiogenesis where the discrepancy between oxygen demand and blood supply is sharpest~\citep{palazon2017hif}.

In summary, the signature genes and functional pathway scores extracted in this section jointly confirm the stroma/immune domain as the primary immunosuppressive hub and the tumour edge as the invasive front transitional zone, and provide a rigorously filtered signature gene set for subsequent cross-sample validation and survival analysis~\citep{zheng2017landscape,schurch2020coordinated}.

\subsection{Spatial Reproducibility of the Niche Gene Programme in the Independent Section CHC20}

This section presents the results of the cross-section consistency validation described in Section 2.7, examining whether the immunosuppressive niche signature gene set extracted from the HCC4R section is generalisable beyond a single tissue section.

\begin{figure}[h]
    \centering
    \begin{subfigure}[b]{0.48\textwidth}
        \centering
        \includegraphics[width=\textwidth]{plots/chc20_immune_module_spatial_score.png}
        \caption{Immune module score.}
        \label{fig:immune_score}
    \end{subfigure}
    \hfill
    \begin{subfigure}[b]{0.48\textwidth}
        \centering
        \includegraphics[width=\textwidth]{plots/chc20_stromal_ecm_module_spatial_score.png}
        \caption{Stromal ECM module score.}
        \label{fig:stromal_score}
    \end{subfigure}
    \vspace{0.4cm}

    \begin{subfigure}[b]{0.48\textwidth}
        \centering
        \includegraphics[width=\textwidth]{plots/chc20_combined_immune_stromal_niche_score.png}
        \caption{Combined microenvironment score.}
        \label{fig:combined_score}
    \end{subfigure}
    \hfill
    \begin{subfigure}[b]{0.48\textwidth}
        \centering
        \includegraphics[width=\textwidth]{plots/hcc4r_vs_chc20_combined_niche_score_distribution.png}
        \caption{Comparison of score distributions.}
        \label{fig:score_comparison}
    \end{subfigure}

    \caption{Spatial distribution and cross-sample comparison of microenvironment module scores.}
    \label{fig:main_microenvironment_figure}
\end{figure}

Figures~\ref{fig:immune_score}, \ref{fig:stromal_score} and \ref{fig:combined_score} display the spatial projections onto CHC20 coordinates of, respectively, the immune module score ($n = 15$ matched genes, encompassing core Treg/TAM molecules including \textit{FOXP3}, \textit{IL2RA}, \textit{CTLA4}, \textit{TIGIT}, \textit{PDCD1}, \textit{LAG3}, \textit{CD163}, \textit{MRC1} and \textit{CXCL12}), the stromal/ECM module score ($n = 23$ matched genes, encompassing CAF/ECM-remodelling molecules including \textit{FAP}, \textit{ACTA2}, \textit{COL1A1}, \textit{COL3A1}, \textit{SPP1}, \textit{MMP2}, \textit{LUM} and \textit{DCN}), and the combined niche score obtained by summing Z-score-standardised values of the two modules. All three maps display highly concordant spatial patterns: high-score regions (yellow/orange) are consistently concentrated in the same band along the upper-right tissue edge, and low-score regions (dark) are consistently concentrated in the large lower area of the section. The dark area corresponds to the tumour-parenchyma-dense region visible in the CHC20 H\&E histological image, demonstrating that signature genes selected from the HCC4R section can effectively identify distinct niches in other sections.

The violin plot in Fig.~\ref{fig:score_comparison} further compares the overall combined niche score distributions between HCC4R and CHC20. The Mann--Whitney U test $p$-value for the two distributions was 0.094, not reaching the conventional 0.05 significance threshold; the medians (HCC4R $\approx -0.5$, CHC20 $\approx -0.2$), interquartile ranges and overall shapes show substantial overlap, and the overall value ranges (HCC4R $\approx -2.9$ to 5.1; CHC20 $\approx -4.5$ to 7.4), while the CHC20 tail is slightly wider, the principal distribution intervals are broadly consistent.

Taken together, the high spatial overlap between CHC20 high-score regions and the histologically defined immunosuppressive zones at the tumour boundary supports the conclusion that the identified niche is a universal histological feature of HCC rather than a section-specific phenomenon~\citep{llovet2021hepatocellular}. It is noteworthy that CHC20 is a combined hepatocellular-cholangiocarcinoma (cHCC-CCA), whose immune microenvironment composition may differ biologically from classical HCC (HCC4R); the score consistency observed between the two sections in this context is therefore particularly meaningful.

\subsection{Independent Prognostic Validation in the TCGA-LIHC Cohort}

\begin{figure}[h]
    \centering
    \begin{subfigure}[b]{0.48\textwidth}
        \centering
        \includegraphics[width=\textwidth]{plots/km_plot.png}
        \caption{Kaplan--Meier curves for ssGSEA high- and low-score groups.}
        \label{fig:km_plot}
    \end{subfigure}
    \hfill
    \begin{subfigure}[b]{0.48\textwidth}
        \centering
        \includegraphics[width=\textwidth]{plots/cox_forest_plot.png}
        \caption{Cox regression forest plot for ssGSEA high- and low-score groups.}
        \label{fig:cox_forest}
    \end{subfigure}
    \caption{Survival analysis results in the TCGA-LIHC cohort.}
    \label{fig:survival_result}
\end{figure}

The HCC4R spatial signature gene set was projected onto TCGA-LIHC patients via ssGSEA, and patients were dichotomised into High and Low groups at the median ssGSEA score. Kaplan--Meier survival analysis (Fig.~\ref{fig:km_plot}) revealed that the two overall survival curves were broadly overlapping, and the log-rank test did not reach statistical significance ($p = 0.79$), indicating that the spatial signature score did not effectively stratify patient overall survival in the current TCGA-LIHC cohort.

Multivariate Cox proportional hazards regression (Fig.~\ref{fig:cox_forest}), incorporating age and tumour stage as additional covariates, yielded a hazard ratio approaching unity for the ssGSEA score (HR = 0.92, 95\% CI: 0.36--2.36, $p = 0.862$); the confidence interval spans 1, indicating that the score does not exhibit independent prognostic value in the current model. In contrast, certain tumour stage variables showed more pronounced trends towards elevated hazard in the Cox model, suggesting that clinical staging remains an important determinant of overall survival in the TCGA-LIHC cohort.

These results indicate that, while the spatial niche signature faithfully reflects tissue cellular composition, it does not function as an independent prognostic factor for bulk overall survival in the TCGA-LIHC cohort, and its relationship with overall survival is limited in this context.

\section{Discussion}

This study established a complete analytical framework spanning five stages---scRNA-seq reference construction, spatial deconvolution, niche identification, mechanistic dissection and cross-sample/clinical validation---achieving the goal of extracting spatial immune regulatory principles from single-cell transcriptomic reference atlases and systematically validating them in independent tissue sections and patient cohorts. Compared with traditional paradigms reliant on prior-defined gene sets, our approach innovatively employs unsupervised Cellular Neighbourhoods clustering to dissect immunosuppressive spatial architecture directly from the underlying data. To overcome the cell dilution effect inherent to the Visium platform, the detection rate difference $\Delta$fraction was introduced to compensate for $\log_2$FC-based filtering, effectively reducing false negatives and improving Treg feature extraction sensitivity. Simultaneous cross-validation of rule-based spatial domain labels against Leiden clustering substantially enhanced the robustness and interpretability of immunosuppressive niche determination. This integrative strategy provides a reproducible methodological paradigm for combining single-cell resolution reference atlases with spatial transcriptomic data; its ``discovery--mechanism--validation'' closed-loop design also has broader implications for spatial subtyping of the TME in other solid tumours.

Despite achieving highly concordant results at the histological level and across independent spatial sections, the TCGA-LIHC independent prognostic validation did not meet expectations: neither Kaplan--Meier analysis ($p = 0.79$) nor multivariate Cox regression (HR = 0.92, $p = 0.862$) reached statistical significance, and neither the forest plot nor the KM curves yielded independent risk-stratification signals from the spatial signature. This discrepancy may be attributable to several methodological limitations. First, projecting ssGSEA onto bulk transcriptomics erases spatial localisation through tissue homogenisation, causing the local high-signal immunosuppressive drivers (e.g., Treg-enriched microregions) to be masked by large-area normal or tumour parenchymal signals~\citep{svensson2018spatiallyvariable,gao2020loss}. In addition, TCGA-LIHC samples are highly heterogeneous in composition, and bulk sequencing is susceptible to sampling-location bias, which may obscure genuine survival associations~\citep{zhang2019landscape}. Furthermore, local niches often drive short-term recurrence or resistance rather than long-term overall survival, making the causal chain to overall survival comparatively long~\citep{zhang2019landscape}; systematic differences in treatment and follow-up protocols across multi-centre cohorts introduce non-biological noise that further impairs prognostic significance~\citep{liu2018tcga}. In summary, while spatial niche signatures accurately and reproducibly characterise the immune spatial architecture, their clinical translation as bulk-level overall survival prognostic factors may require more refined outcome endpoints (e.g., recurrence-free survival, RFS) or more uniformly collected prospective samples~\citep{liu2018tcga}.

Beyond the non-significant survival analysis, several additional methodological limitations remain to be addressed. The immunosuppressive niche score depends moderately on the choice of neighbourhood parameters ($k$ and radius multiplier); although sensitivity analyses have confirmed overall robustness ($\rho > 0.90$, WJ $> 0.80$), residual local spatial boundary effects under extreme parameter changes cannot be entirely dismissed. Future work could adopt Bayesian hierarchical models that simultaneously perform probabilistic inference on both posterior spatial patterns and neighbourhood probability distributions, reducing reliance on a single fixed neighbourhood definition~\citep{kleshchevnikov2022cell2location}. The ligand--receptor communication analysis in Section 2.6, while employing rank normalisation to circumvent the zero-inflation artefact of raw product scores, cannot fully distinguish genuine bidirectional communication from spurious high co-expression arising from different cell types each individually expressing ligand or receptor genes within the same spot; addressing this may require coupling with probabilistic models of ligand--receptor gene co-localisation. The immunosuppressive gene module was selected based solely on upregulation, without considering downregulation of inhibitory molecules (e.g., HLA gene loss) as an additional immune evasion mechanism~\citep{montesion2021hla}. Moreover, the current analysis incorporated only one HCC section (HCC4R) and one combined hepatocellular-cholangiocarcinoma section (CHC20) for cross-sample validation; sample size remains insufficient, and future work should extend to multi-centre, multi-cancer-type spatial transcriptomic cohorts.

Looking ahead, the analytical framework established in this study can be extended and applied in several directions. First, integration with multiplexed immunofluorescence (mIF) or spatial proteomics platforms such as MIBI/CODEX would enable protein-level validation of the spatial co-localisation of key molecules in the immunosuppressive niche (CXCL12--CXCR4 axis, PD-1/PD-L1 immune checkpoints), providing molecular biomarkers for personalised prediction of immune checkpoint blockade (ICB) treatment response~\citep{keren2018structured,zhu2021spark}. As higher-resolution spatial sequencing technologies---including 10x Visium HD, Slide-seq and Stereo-seq---continue to mature, it will be possible to re-evaluate whether the domain definitions and signature genes established in this study remain applicable at sub-spot spatial resolution approaching single-cell or sub-cellular precision, further exploring the fine topology of the tumour--immune interface~\citep{moses2022museum,chen2022spatia}. In addition, the ``spatial niche signature'' could be used as an input feature to train Transformer-based deep-learning models for marker-free prediction of recurrence risk at surgical margins, or integrated with digital pathology to combine macroscopic H\&E image features with microscopic spatial transcriptomic signals in a multimodal framework for cross-scale tumour immune status assessment~\citep{chen2016liver,shmatko2022artific}. The data-processing pipelines, feature-selection criteria and sensitivity-analysis frameworks accumulated in this study can serve as a foundational methodological reference for all of the above future directions.

% ============================================================
%  Appendices
% ============================================================

\newpage
\appendix
\renewcommand{\thesection}{\Alph{section}}

\section{Ligand--Receptor Communication Pair List}
\label{app:lr_table}

\begin{table}[H]
\centering
\caption{Ligand--receptor communication pairs included in the analysis.}
\label{tab:lr_pairs}
\begin{tabularx}{\textwidth}{lllX}
\toprule
\textbf{Ligand} & \textbf{Receptor} & \textbf{Pair} & \textbf{Biological significance} \\
\midrule
CXCL12 & CXCR4   & CXCL12--CXCR4  & CAF-secreted; recruits Tregs and MDSCs \\
TGFB1  & TGFBR1  & TGFB1--TGFBR1  & Immunosuppressive cytokine; co-secreted by CAFs and Tregs \\
PDCD1  & CD274   & PD-1--PD-L1    & Classical immune checkpoint axis \\
TIGIT  & NECTIN2 & TIGIT--NECTIN2 & Treg inhibitory checkpoint \\
LAG3   & HLA-DRA & LAG3--MHC-II   & T-cell exhaustion-associated \\
CCL22  & CCR4    & CCL22--CCR4    & TAM-mediated Treg recruitment \\
IL10   & IL10RA  & IL10--IL10RA   & Treg immunosuppressive cytokine axis \\
SPP1   & CD44    & SPP1--CD44     & TAM-secreted osteopontin \\
MIF    & CD74    & MIF--CD74      & Macrophage migration inhibitory factor \\
VEGFA  & KDR     & VEGFA--KDR     & Angiogenesis; CAF-associated \\
LGALS9 & HAVCR2  & Galectin9--TIM-3 & Inhibitory checkpoint \\
CCL2   & CCR2    & CCL2--CCR2     & Macrophage recruitment pathway \\
\bottomrule
\end{tabularx}
\end{table}

\section{Definition of Neighbourhood Parameter Sensitivity Analysis Metrics}
\label{app:sensitivity}

The neighbourhood parameter sensitivity analysis described in Section 2.4 employs two complementary continuous metrics to jointly assess the consistency of composite niche score distributions before and after perturbation of $k$-NN neighbour count and radius multiplier.

The first metric is Spearman's rank correlation coefficient $\rho$, defined as
\begin{equation}
\rho = 1 - \frac{6 \sum_i d_i^2}{n(n^2-1)}
\label{eq:spearman}
\end{equation}
where $d_i$ is the difference in ranks between the reference and perturbed scores for spot $i$, and $n$ is the total number of spots. $\rho$ directly measures the consistency of the relative ordering of all spots and is unaffected by any cutoff threshold, with values in $[-1, 1]$; values approaching 1 indicate that the parameter perturbation has virtually no effect on spot ranking.

The second metric is the continuous weighted Jaccard similarity, defined as
\begin{equation}
\text{WJ}(\mathbf{a}, \mathbf{b}) = \frac{\sum_i \min(a_i, b_i)}{\sum_i \max(a_i, b_i)}
\label{eq:weighted_jaccard}
\end{equation}
where $\mathbf{a}$ and $\mathbf{b}$ are the reference and perturbed score vectors after Min-Max normalisation to $[0,1]$. Compared with the traditional binary set-based Jaccard index, this continuous formulation naturally assigns higher weight to high-scoring regions (i.e., candidate niche regions), compensating for the limitation of Spearman's rank correlation in treating top and bottom ranks equally. Values in $[0, 1]$ approaching 1 indicate high spatial overlap of high-score regions.

\section{Auxiliary Score Formulae}
\label{app:scores}

The auxiliary scoring system is defined as follows, where $z(\cdot)$ denotes Z-score standardisation across spots:

\begin{align}
S_{\text{Treg-like},i} &= z(\hat{p}_{\text{Treg},i}) + z(\hat{s}_{\text{gene},i}) \\
S_{\text{immune-stroma},i} &= z(\hat{p}_{\text{Treg},i}) + z(\hat{p}_{\text{T/NK},i}) + z(\hat{p}_{\text{Myeloid},i}) + z(\hat{p}_{\text{Fibroblast},i}) \\
S_{\text{niche},i} &= z(\hat{p}_{\text{Treg},i}) + z(\hat{p}_{\text{Myeloid},i}) + z(\hat{p}_{\text{Fibroblast},i}) + z(\hat{s}_{\text{gene},i}) + z(\bar{n}_{\text{Hep},i})
\end{align}

where $\hat{p}$ denotes the Cell2location posterior abundance normalised to proportions, $\hat{s}_{\text{gene}}$ denotes the immunosuppressive gene module score, and $\bar{n}_{\text{Hep}}$ denotes the mean hepatocyte proportion within the spot neighbourhood (used to capture tumour--stroma interface effects). Dual-module scores (used for CHC20 validation):

\begin{align}
\text{immune\_score}_i &= \frac{1}{|G_{\text{immune}}|}\sum_{g \in G_{\text{immune}}} E_{ig} \\
\text{stromal\_score}_i &= \frac{1}{|G_{\text{stromal}}|}\sum_{g \in G_{\text{stromal}}} E_{ig} \\
\text{combined\_niche\_score}_i &= z(\text{immune\_score}_i) + z(\text{stromal\_score}_i)
\end{align}

$E_{ig}$ denotes log1p-normalised (CP10K) expression; the immune module comprises 21 genes including \textit{FOXP3}, \textit{IL2RA}, \textit{CTLA4} and \textit{CD163}, and the stromal/ECM module comprises 21 genes including \textit{FAP}, \textit{ACTA2}, \textit{COL1A1} and \textit{SPP1}.

\section{Core Code Implementation}
\label{app:code}

\subsection{Cell2location Deconvolution}

\begin{lstlisting}[caption={Core code for RegressionModel training and spatial modelling}]
# Stage 1: Train reference signature matrix
from cell2location.models import RegressionModel
RegressionModel.setup_anndata(adata_sc, labels_key='final_celltype')
mod = RegressionModel(adata_sc)
mod.train(max_epochs=250, batch_size=1024, use_gpu=True,
          early_stopping=True, early_stopping_patience=30)
cell_state_df = mod.export_posterior(adata_sc,
    use_cols=mod.feature_in_columns)['means_per_cluster_mu_fg']

# Stage 2: Spatial modelling
import cell2location
cell2location.models.Cell2location.setup_anndata(adata_vis)
mod = cell2location.models.Cell2location(adata_vis,
    cell_state_df=cell_state_df,
    N_cells_per_location=30, detection_alpha=20)
mod.train(max_epochs=1000, batch_size=1000, use_gpu=True)
adata_vis = mod.export_posterior(adata_vis,
    sample_kwargs={'num_samples': 1000, 'batch_size': 1000})
# Results stored in adata_vis.obsm['means_cell_abundance_w_sf']
\end{lstlisting}

\subsection{Neighbourhood Composition Clustering and Niche Annotation}

\begin{lstlisting}[caption={Neighbourhood composition vector computation, Leiden clustering and composite rank annotation}]
from scipy.spatial import KDTree
import scanpy as sc

# Build kNN adjacency graph and compute neighbourhood composition vectors
tree = KDTree(coords)
_, indices = tree.query(coords, k=k + 1)  # Exclude self
neighbors = [indices[i, 1:] for i in range(len(coords))]
arr = proportions.to_numpy()
neighborhood_comp = np.array(
    [arr[idx].mean(axis=0) if len(idx) else arr[i]
     for i, idx in enumerate(neighbors)])

# Leiden clustering (directly in composition vector feature space, no PCA)
adata_nc = ad.AnnData(neighborhood_comp.astype(float))
sc.pp.neighbors(adata_nc, n_neighbors=n_neighbors, use_rep='X')
sc.tl.leiden(adata_nc, resolution=0.5, key_added='nc_cluster')

# Composite rank annotation of immunosuppressive niche
stats = df.groupby('nc_cluster')[
    [treg_col, myeloid_col, fibroblast_col, gene_score_col]].mean()
for col in stats.columns:
    stats[f'rank_{col}'] = stats[col].rank(ascending=False)
stats['combined_rank'] = stats.filter(like='rank_').sum(axis=1)
niche_id = stats['combined_rank'].idxmin()
\end{lstlisting}

\subsection{Two-Layer Signature Gene Filtering}

\begin{lstlisting}[caption={Layer 1 differential testing and AUC calculation}]
# Layer 1: block-level Mann-Whitney U + AUC
from scipy.stats import mannwhitneyu
# Spatial grid aggregation (reduces spatial autocorrelation)
gx = np.floor(df['spatial_x'].values / block_size).astype(int)
gy = np.floor(df['spatial_y'].values / block_size).astype(int)
block_ids = [f'{x}_{y}' for x, y in zip(gx, gy)]
# Compute AUC = U / (n_high * n_low)
u_stat, pval = mannwhitneyu(high_vals, low_vals, alternative='greater')
auc = u_stat / (len(high_vals) * len(low_vals))

# Gini Index for focal/sparse gene filtering
def gini_index(x):
    x = np.sort(x.ravel())
    n = len(x)
    cumx = np.cumsum(x)
    return (n + 1 - 2 * cumx.sum() / cumx[-1]) / n if cumx[-1] > 0 else 0
\end{lstlisting}

\subsection{Dual-Module Scoring and ssGSEA}

\begin{lstlisting}[caption={Dual-module score computation (Python) and ssGSEA scoring (R)}]
# Python: dual-module scoring
def compute_dual_module_scores(adata, immune_genes, stromal_genes):
    expr = _expression_frame(adata)  # CP10K + log1p
    immune_score = expr[[g for g in immune_genes
                          if g in expr.columns]].mean(axis=1)
    stromal_score = expr[[g for g in stromal_genes
                           if g in expr.columns]].mean(axis=1)
    def _z(s): return (s - s.mean()) / s.std() if s.std() > 0 else s * 0
    return immune_score, stromal_score, _z(immune_score) + _z(stromal_score)
\end{lstlisting}

\begin{lstlisting}[language=R, caption={ssGSEA and multivariate Cox regression (R)}]
library(GSVA); library(survival); library(survminer)
gsva_scores <- gsva(expr_matrix, list(sig = signature_genes),
    method='ssgsea', kcdf='Gaussian', abs.ranking=TRUE, min.sz=5)
cox_model <- coxph(
    Surv(OS_time, OS_event) ~ ssgsea_score + age + stage,
    data = clinical_data)
summary(cox_model)
\end{lstlisting}

% ============================================================
% References
% ============================================================
\bibliographystyle{unsrt}
\bibliography{report_0701}

\end{document}
